\documentclass{beamer}
\usepackage[T2A]{fontenc}
\usepackage[utf8]{inputenc}
\usepackage[russian]{babel}
\usepackage{amsmath, amssymb, amsthm}
\usepackage{amsfonts}
\usepackage{graphicx}
\usepackage{booktabs}

% Настройка темы оформления
\usetheme{Madrid}
\usecolortheme{default}

% Настройка заголовка
\title{Задача о минимальной норме интерполяционного проектора}
\author{Зайков Иван Викторович}
\institute{
    Ярославский государственный университет им. П. Г. Демидова\\
    Кафедра математического анализа
    \and
    Группа: ПМИ-21МО
    \and
    Научный руководитель: к.ф.-м.н., доцент Ухалов Алексей Юрьевич
}
\date{2026}

\begin{document}
\setbeamertemplate{footline}[frame number]

% Титульный слайд
\begin{frame}
    \titlepage
\end{frame}

\begin{frame}{Постановка задачи}
    \textbf{Интерполяция:} для функции $f(x)$ на кубе $Q_n = [0,1]^n$ строим линейный полином
    \[
    p(x) = a_1 x_1 + \dots + a_n x_n + b,
    \]
    совпадающий с $f$ в $n+1$ узле (вершинах симплекса $S$).

    \vspace{0.3cm}

    \textbf{Интерполяционный проектор} $P: C(Q_n) \to \Pi_1(\mathbb{R}^n)$

    \vspace{0.3cm}

    \textbf{Норма проектора (константа Лебега):}
    \[
    \|P\| = \max_{x \in \operatorname{ver}(Q_n)} \sum_{j=1}^{n+1} |\lambda_j(x)|,
    \]
    где $\lambda_j$ — базисные многочлены Лагранжа (барицентрические координаты).

    \vspace{0.3cm}

    \textbf{Цель:} используя симплекс максимального объёма, построенный из матрицы Адамара, получить верхнюю оценку минимальной нормы проектора $\theta_{31}$ в $\mathbb{R}^{31}$.
\end{frame}

% === 2. Используемые обозначения ===
\begin{frame}{Обозначения}
    \begin{itemize}
        \item $Q_n = [0,1]^n$ — единичный куб
        \item $C(Q_n)$ — пространство непрерывных функций с нормой $\|f\| = \max_{x \in Q_n} |f(x)|$
        \item $\Pi_1(\mathbb{R}^n) = \operatorname{span}\{1, x_1, \dots, x_n\}$ — пространство полиномов от $n$ переменных степени $\le 1$
        \item $S \subset Q_n$ — невырожденный симплекс (вершины $x^{(1)}, \dots, x^{(n+1)}$)
    \end{itemize}
\end{frame}

% === 3. Интерполяционный проектор ===

\subsection{Постановка задачи интерполяции} % (в презентации заголовок не нужен, просто для логики)

\subsection{Базисные многочлены Лагранжа и барицентрические координаты}
\begin{frame}{Базисные многочлены Лагранжа}
    \textbf{Матрица $A$} (строки — вершины, последний столбец — единицы):
    \[
    A = \begin{pmatrix}
    x_1^{(1)} & \dots & x_n^{(1)} & 1 \\
    \vdots & & \vdots & \vdots \\
    x_1^{(n+1)} & \dots & x_n^{(n+1)} & 1
    \end{pmatrix}, \qquad \det A \neq 0
    \]

    \vspace{0.3cm}

    $\Delta = \det(A)$, $\Delta_j(x)$ — замена $j$-й строки на $(x_1,\dots,x_n,1)$.

    \vspace{0.3cm}

    \textbf{Базисные многочлены Лагранжа:}
    \[
    \lambda_j(x) = \frac{\Delta_j(x)}{\Delta}
    \]

    \vspace{0.3cm}

    \textbf{Через обратную матрицу:} $\lambda(x) = A^{-1} b(x)$, $b(x) = (x_1,\dots,x_n,1)^T$.

    \vspace{0.3cm}

    \textbf{Свойства:} $\lambda_j(x^{(k)}) = \delta_{jk}$, $\sum \lambda_j(x) = 1$.
\end{frame}
\subsection{Интерполяционный проектор}
\begin{frame}{Интерполяционный проектор}
    \textbf{Определение:} $P: C(Q_n) \to \Pi_1(\mathbb{R}^n)$, $P f(x^{(j)}) = f(x^{(j)})$, $j=1,\dots,n+1$.

    \vspace{0.2cm}

    \textbf{Формула Лагранжа:}
    \[
    P f(x) = \sum_{j=1}^{n+1} f(x^{(j)}) \lambda_j(x)
    \]

    \vspace{0.2cm}

    \textbf{Норма:} $\displaystyle \|P\| = \max_{x \in \operatorname{ver}(Q_n)} \sum_{j=1}^{n+1} |\lambda_j(x)|$

    \vspace{0.2cm}

    \textbf{Минимальная норма:} $\theta_n = \min \|P\|$
\end{frame}
\subsection{Неравенство Лебега}  % (новый слайд)
\begin{frame}{Неравенство Лебега}
    Пусть $E_n(f)$ — наилучшее приближение функции $f$ полиномами из $\Pi_1(\mathbb{R}^n)$:
    \[
    E_n(f) = \min_{p \in \Pi_1(\mathbb{R}^n)} \|f - p\|_{C(Q_n)}.
    \]

    \vspace{0.3cm}

    \textbf{Неравенство Лебега:}
    \[
    \|f - P f\|_{C(Q_n)} \le (1 + \|P\|) \cdot E_n(f).
    \]

    \vspace{0.3cm}

    \textbf{Следствие:} чем меньше норма проектора $\|P\|$, тем ближе интерполяционный полином к наилучшему приближению.
\end{frame}

\subsection{Оценки нормы проектора}
\begin{frame}{Общие оценки $\theta_n$}
    \textbf{Неравенства:}
    \[
    \frac{1}{4}\sqrt{n} < \theta_n < 3\sqrt{n}, \qquad 3 - \frac{4}{n+1} \le \theta_n
    \]

    \vspace{0.3cm}

    \textbf{Асимптотическая эквивалентность:}
    \[
    \|P\| \asymp \theta_n \quad\Leftrightarrow\quad c_1\sqrt{n} \le \theta_n \le c_2\sqrt{n}
    \]

    \vspace{0.3cm}

    \textbf{Идея:} использовать симплекс максимального объёма для верхней оценки $\theta_n$.
\end{frame}

% === 4. Матрицы Адамара ===
\subsection{Определение и основные свойства}
\begin{frame}{Матрица Адамара}
    \textbf{Определение:} $H$ — квадратная \textbf{$(-1/1)$-матрица} (элементы $\pm 1$), удовлетворяющая условию:
    \[
    H H^T = n I_n
    \]

    \vspace{0.3cm}

    \textbf{Свойства:}
    \begin{itemize}
        \item $|\det H| = n^{n/2}$ — \textbf{максимальный определитель}
        \item Существует при $n = 1, 2, 4k$ (гипотеза Адамара: для любого $n = 4k$)
        \item Для порядка $32$: $13\,710\,027$ неэквивалентных матриц
    \end{itemize}

    \vspace{0.3cm}

    \textbf{Построение (Сильвестр):}
    \[
    H_{1} = (1), \quad H_{2n} = \begin{pmatrix} H_n & H_n \\ H_n & -H_n \end{pmatrix}
    \]

    \vspace{0.3cm}

    \textbf{Связь с симплексом:}
    Матрица Адамара $H$ порядка $n+1$ даёт максимальный определитель $\det H = (n+1)^{(n+1)/2}$, который после преобразования становится максимальным $(0/1)$-определителем порядка $n$.
\end{frame}

\subsection{Методы построения матриц Адамара}
% Метод Сильвестра (уже есть в слайде 6, можно оставить как есть или дублировать)

\subsubsection{Способы построения матриц Адамара}
\begin{frame}{Типы матриц Адамара порядка 32}
    \begin{itemize}
        \item \texttt{syl} — метод Сильвестра
        \item \texttt{pal} — метод Пэли (на основе конечных полей)
        \item \texttt{t1}, \texttt{t2}, \texttt{t3}, \texttt{t4} — метод тензорного произведения
    \end{itemize}

    \vspace{0.3cm}

    Шесть матриц, построенных указанными способами, взяты из библиотеки Нила Слоуна (доступна по адресу \url{https://neilsloane.com/hadamard/}).
\end{frame}


\subsection{Количество неэквивалентных матриц Адамара} % (новый слайд)
\begin{frame}{Количество матриц Адамара}
    \begin{table}
    \centering
    \begin{tabular}{c|c}
    \hline
    Порядок $n$ & Количество неэквивалентных матриц \\
    \hline
    1, 2, 4, 8, 12 & 1 \\
    16 & 5 \\
    20 & 3 \\
    24 & 60 \\
    28 & 487 \\
    32 & $13\,710\,027$ \\
    \hline
    \end{tabular}
    \end{table}

    \vspace{0.3cm}

    Ввиду огромного числа матриц Адамара порядка $32$ (более $13$ миллионов) полный перебор всех возможных симплексов не проводился. Исследование ограничено шестью представителями вышеуказанных типов.
\end{frame}

\subsection{Переход к (0/1)-матрицам и построение симплекса}
\begin{frame}{Построение симплекса}
    \textbf{Исходные данные:} матрица Адамара $H$ порядка $32$ (элементы $\pm 1$).

    \vspace{0.3cm}

    \textbf{Шаги:}
    \begin{enumerate}
        \item Приводим $H$ к виду $B$, где первый столбец и первая строка состоят из единиц (умножением строк/столбцов на $-1$).
        \item Удаляем первую строку и первый столбец → матрица $C$ порядка $31$.
        \item Заменяем элементы: $-1 \to 1$, $1 \to 0$ → получаем $(0/1)$-матрицу $D$ порядка $31$.
        \item Добавляем нулевую строку (вершину) → $32$ вершины в $\mathbb{R}^{31}$.
    \end{enumerate}

    \vspace{0.3cm}

    \textbf{Результат:} $V = \{0, \operatorname{row}_1(D), \dots, \operatorname{row}_{31}(D)\} \subset [0,1]^{31}$ — симплекс максимального объёма.
\end{frame}

% === 5. Описание программной реализации ===
\subsection{Общая архитектура} % (в презентации слайд "Программная реализация" должен объединять несколько подразделов)
\begin{frame}{Программная реализация (1/2)}
    \textbf{Основные классы:}
    \begin{itemize}
        \item `ISimplexProvider` — интерфейс для получения вершин симплекса
        \item `SilvesterMethodProvider`, `HadamardFileProvider` — реализация
        \item `FullEnumUpperBoundCalculator` — многопоточный перебор вершин
        \item `LebesgueFunction` — вычисление суммы модулей
        \item `HadamardMatrixIterator` — чтение матриц из файла
    \end{itemize}
    \textbf{Технологии:} C++17, Eigen, CMake, многопоточность (атомарные операции).
\end{frame}

\subsection{Ввод и вывод данных}
\begin{frame}{Программная реализация (2/2)}
    \textbf{Режимы работы:}
    \begin{itemize}
        \item Построение симплекса по Сильвестру
        \item Загрузка матриц Адамара из файла
    \end{itemize}

    \textbf{Параметры командной строки:}
    \begin{itemize}
        \item \texttt{--dimension N} — размерность пространства $\mathbb{R}^N$
        \item \texttt{--load-hadamard FILE} — загрузить матрицу Адамара из файла
        \item \texttt{--threads N} — количество потоков
        \item \texttt{--range START END} — частичный перебор вершин в диапазоне $[\text{START}, \text{END})$ (для распределения вычислений)
    \end{itemize}
\end{frame}

\subsection{Вычисление нормы проектора}
\begin{frame}{Полный перебор}
    \textbf{Число вершин гиперкуба $Q_{31}$:}
    \[
    N = 2^{31} = 2\,147\,483\,648
    \]

    \vspace{0.3cm}

    \textbf{Особенности реализации:}
    \begin{itemize}
        \item Кодирование вершины 31-битовым целым числом
        \item Многопоточность (16 потоков)
        \item Атомарные операции для синхронизации
        \item Логирование каждые $10^7$ вершин
    \end{itemize}

    \vspace{0.3cm}

    \[
    \theta_{31} = \max_{x \in \{0,1\}^{31}} \sum_{j=1}^{32} |\lambda_j(x)|
    \]
\end{frame}

% === 6. Результаты вычисления ===
\subsection{Параметры вычислительного эксперимента}
\begin{frame}{Параметры вычислительной системы}
    \textbf{Аппаратное обеспечение:}
    \begin{table}
    \centering
    \begin{tabular}{l|c}
    \hline
    \textbf{Параметр} & \textbf{Значение} \\
    \hline
    Процессор & Intel Core i5-1240P (12-го поколения) \\
    Ядра (потоки) & 16 \\
    Оперативная память & 15 ГБ \\
    \hline
    \end{tabular}
    \end{table}

    \vspace{0.5cm}

    \textbf{Программное обеспечение:}
    \begin{table}
    \centering
    \begin{tabular}{l|c}
    \hline
    \textbf{Параметр} & \textbf{Значение} \\
    \hline
    Операционная система & Ubuntu 24.04.3 LTS \\
    Компилятор & GCC 13.3.0 \\
    Библиотека & Eigen 3 \\
    Количество потоков & 16 \\
    \hline
    \end{tabular}
    \end{table}
\end{frame}

\subsection{Результаты и сравнение с теоретической оценкой}
\begin{frame}{Результаты}
    \textbf{Полученное значение:}
    \[
    \boxed{\theta_{31} \le 5.000}
    \]

    \vspace{0.5cm}

    \begin{table}
    \centering
    \begin{tabular}{l|c}
    \hline
    \textbf{Параметр} & \textbf{Значение} \\
    \hline
    Теоретическая оценка & $\sqrt{32} \approx 5.656854$ \\
    Экспериментальное значение & $5.000$ \\
    Время выполнения & $290$ секунд ($\approx 5$ мин) \\
    Скорость обработки & $\approx 7.4 \times 10^6$ вершин/сек \\
    \hline
    \end{tabular}
    \end{table}

    \vspace{0.3cm}
\end{frame}

% === Заключение ===
\begin{frame}{Заключение}
    \textbf{Выполнено:}
    \begin{enumerate}
        \item Построен симплекс максимального объёма в $\mathbb{R}^{31}$ на основе матрицы Адамара
        \item Реализован многопоточный полный перебор $2^{31}$ вершин гиперкуба
        \item Получена верхняя оценка $\theta_{31} \le 5.000$
    \end{enumerate}

    \vspace{0.5cm}
\end{frame}

% === Спасибо за внимание ===
\begin{frame}
    \centering
    \Huge \textcolor{blue}{Спасибо за внимание!}

    \vspace{0.5cm}
\end{frame}
\end{document}